\chapter{Navier--Stokes: terms to put back}

\textbf{Purpose:} Collect every known member of $P_{\mathrm{all}}$ and $F_{\mathrm{all}}$ that
equation (1) drops, with the term that puts it back and the source that shows it present; name what
is not known; and look across the known members for a pattern left over. \textbf{Scope:} the
research paper chapters on the dropped terms, the average and the two sets, whose propositions this
chapter names by number, and the sources registered with them.

\section{The method}

The approach to the properties of a fluid has three steps. Examine everything known. Examine what
is not known, and name it as not known. Then look across the parts that are known and see whether
any pattern is left. The three sections below are those steps in order.

$P_{\mathrm{all}}$ is the set of all properties of a fluid and $F_{\mathrm{all}}$ the set of all
forces a fluid can experience, known and not yet known, and neither list is claimed complete.
Equation (1) keeps $P_{(1)} = \{\nu\}$ and $F_{(1)} = \{-\nabla p,\ \nu\Delta u,\ f\}$.

\section{What is known}

Each row is a member shown present in a real fluid and dropped by equation (1). The status column
uses the status names of the earlier chapters, and \emph{proved here} marks a proof short enough to be
written in the row itself.

\begingroup
\small
\begin{longtable}{@{}>{\raggedright\arraybackslash}p{0.20\textwidth}%
>{\raggedright\arraybackslash}p{0.19\textwidth}%
>{\raggedright\arraybackslash}p{0.25\textwidth}%
>{\raggedright\arraybackslash}p{0.28\textwidth}@{}}
\toprule
Member & What equation (1) does & Term to put back & What shows it present \\
\midrule
\endhead
\bottomrule
\endfoot
\multicolumn{4}{@{}l}{\textbf{In} $P_{\mathrm{all}}$} \\[0.3em]
Viscosity that varies with temperature
& holds $\nu$ constant
& $2\mu'(T)\,D\nabla T$ in the momentum balance; Propositions 1 and 3
& computed: $-2.45$ percent per kelvin at $20\,^\circ$C, IAPWS R12-08~\cite{src:IAPWS-R12-08} \\[0.5em]
Viscosity that varies with pressure
& holds $\nu$ constant
& $2(\partial\mu/\partial p)\,D\nabla p$
& computed: $-0.031$ percent per megapascal at $20\,^\circ$C, IAPWS R12-08~\cite{src:IAPWS-R12-08}; the sign turns
positive between $25$ and $40\,^\circ$C \\[0.5em]
Density that varies with temperature
& sets $\rho = 1$
& $\rho(T)$ in the inertia, and $\rho'(T)\,DT/Dt + \rho\theta = 0$; Proposition 4
& computed: expansion of $2.07\times 10^{-4}$ per kelvin at $20\,^\circ$C, zero at
$3.978\,^\circ$C, IAPWS R6-95 (2018)~\cite{src:IAPWS-R6-95-2018} \\[0.5em]
Density that varies with pressure: the weight of the fluid above, and stirring
& sets $\rho = 1$
& $\rho(p,T)$ in the inertia and the mass balance, with the compressibility
$\partial\ln\rho/\partial p$
& computed: $4.589\times 10^{-10}$ per pascal at $20\,^\circ$C, a column at rest denser at the
bottom by $4.49\times 10^{-6}$ per meter, and a turning cup less dense at the bottom center and
denser at the bottom wall; IAPWS R6-95 (2018)~\cite{src:IAPWS-R6-95-2018}, intuition I6 \\[0.5em]
A finite speed of sound
& condition (2) makes it infinite; Proposition 12
& $\theta \neq 0$ and an equation of state $p(\rho,T)$
& computed: $1482.35$ meters per second at $20\,^\circ$C, IAPWS R6-95 (2018)~\cite{src:IAPWS-R6-95-2018} \\[0.5em]
Bulk viscosity
& removed with $\theta$ by condition (2)
& $(\lambda+\mu)\nabla\theta + \theta\nabla\lambda$, with $\lambda + \tfrac23\mu$ the bulk viscosity
& read: about three times the shear viscosity from 7 to $50\,^\circ$C, $2.469\times 10^{-3}$
Pa\,s at $25\,^\circ$C; Holmes, Parker and Povey~\cite{src:Holmes-Parker-Povey-2011} \\[0.5em]
Heat: a temperature field, its conduction, and the heat shear puts in
& holds no temperature
& $\rho c\,DT/Dt = \nabla\cdot(k\nabla T) + 2\mu|D|^2$; Propositions 2 and 5
& computed: $k = 0.5980$ watts per meter per kelvin at $20\,^\circ$C, IAPWS R15-11~\cite{src:IAPWS-R15-11} \\[0.5em]
A level of pressure at which the liquid breaks
& holds no level of pressure; Proposition 13
& pressure fixed in level, and an equation of state that holds the phase boundary as a free
surface; the coupled system of the research paper, Propositions 14 to 17
& read: $-168$ MPa from nucleation theory at $20\,^\circ$C, about $-140$ MPa in inclusions, a few
to a few tens of megapascals below zero by other methods; Caupin and Herbert~\cite{src:Caupin-and-Herbert-2006} \\[0.5em]
Molecules, and the field between them
& averages both over a cell; Propositions 8 to 11
& none written yet
& proved, and computed: $3.337\times 10^{28}$ molecules per cubic meter at $20\,^\circ$C, IAPWS
R6-95 (2018)~\cite{src:IAPWS-R6-95-2018} \\[0.8em]
\multicolumn{4}{@{}l}{\textbf{In} $F_{\mathrm{all}}$} \\[0.3em]
Thermal fluctuating stress
& holds none
& $\nabla\cdot\tilde\tau$, Gaussian with covariance $\tfrac{2\nu k_BT}{\rho}
(\delta_{ik}\delta_{jl}+\delta_{il}\delta_{jk}-\tfrac23\delta_{ij}\delta_{kl})\,
\delta^3(x-x')\,\delta(t-t')$
& read: the scale at which it matters is about the Kolmogorov length; Eyink, Bandak, Goldenfeld and
Mailybaev~\cite{src:Eyink-Bandak-Goldenfeld-Mailybaev-2022} \\[0.5em]
Light, from energy concentrated in a collapse
& holds no channel for it
& a loss of energy to radiation in the balance of energy
& read: flashes shorter than $50$ ps from collapsing bubbles, energy concentrated by twelve orders of
magnitude, brightness up a hundredfold from $40$ to $0\,^\circ$C; Barber, Hiller, Lofstedt,
Putterman and Weninger~\cite{src:Barber-1997}; intuition I11 \\[0.5em]
Gravity acting on a density that varies
& takes $f$ from outside, and $f = 0$ in statements (A) and (B)
& $\rho(T)\,g$
& proved here: $\nabla\times(\rho g) = \nabla\rho\times g$, nonzero wherever $\nabla T$ is not
parallel to $g$ and $\rho'(T)\neq 0$, and a force with nonzero curl cannot be taken into
$-\nabla p$ \\[0.5em]
The force of statements (C) and (D)
& admits any smooth $f$ with the decay of condition (5)
& none: the want is to name the member of $F_{\mathrm{all}}$ that $f$ stands for
& open \\
\end{longtable}
\endgroup

\section{What is not known}

\textbf{Members not yet measured.} By the definition of the two sets, they cannot be listed. Their
place in this chapter is the statement that the table above is not the whole of either set.

\textbf{Intuitions not yet tested.} Each is held in the author's words, with the source that would
test it and what would count as a hit or a miss, in the log of intuitions recorded before reading
that sits beside this workbook. An entry there is scored only after the source is read.

\textbf{Candidates with no held source.} Memory in the stress, gas dissolved in the liquid, and
forces at a free surface are candidates. None is entered until a source is held and read.

\textbf{Questions the earlier chapters leave open.} Whether the viscosity of water depends on $|D|$
at rates a solution reaches. Whether a flow can hold its own temperature gradient inside the null
space of its own rate of strain. The time form of Proposition 10. Whether spherical harmonics of
the motion recover the container. Whether keeping any member above changes whether solutions break
down: that is open for every row.

\textbf{An analogy, held as theory.} Microcavitation is to water what quantum foam is to space: a
fluid under so much tension that it must move angular momentum between the bulk and a boundary,
and it can do so only through heat. The tension is in the fluid itself and not in a surface.

Two parts of this are read. Caupin and Herbert~\cite{src:Caupin-and-Herbert-2006} describe a liquid stretched below its saturated
vapor pressure, with a pressure that can go negative, and read a negative pressure as a density
reduced from the equilibrium one: the fluid itself is pulled apart. They describe homogeneous
nucleation as the birth of the new phase by thermal fluctuations, at a rate proportional to
$\exp[-E_b/(k_BT)]$, over a barrier $E_b = 16\pi\sigma^3/\bigl(3(P_\sigma - P)^2\bigr)$. The tension
of the fluid is the factor $P_\sigma - P$, and the more of it the lower the barrier. The surface
tension $\sigma$ enters only as the cost of the wall of the bubble once one forms, and it is a
further member of $P_{\mathrm{all}}$ that equation (1) does not hold. The break is crossed by heat.

The analogy goes further. Introducing space into the fluid sets off a chain, the energy available
compounds for a short while and grows exponentially, and the end of the chain is the
microcavitation. With $\Delta P = P_\sigma - P$ the tension and $A = 16\pi\sigma^3/(3k_BT)$, the
links are these.

\begin{enumerate}
\item Space is introduced: the fluid is stretched and its density falls below the equilibrium one.
\emph{Read}, Caupin and Herbert.
\item The pressure falls below the saturated vapor pressure and can go negative: the fluid is under
tension $\Delta P$. \emph{Read}, the same.
\item The barrier to a bubble falls as the tension rises, $E_b = A\,k_BT/\Delta P^2$, as the inverse
square of the tension. \emph{Read}, their equation (2).
\item The rate of nucleation, proportional to $\exp(-A/\Delta P^2)$, grows exponentially:
\[
\frac{d}{d\,\Delta P}\,\ln(\text{rate}) = \frac{2A}{\Delta P^3} > 0 ,
\]
and climbs from negligible to certain over a narrow range of tension. \emph{Proved here} from
their equation (2).
\item Heat carries one fluctuation over the barrier: a bubble reaches the critical radius
$R_c = 2\sigma/\Delta P$. \emph{Read}: nucleation is the birth of the new phase by thermal
fluctuations.
\item Past $R_c$ the work to form a bubble, $-\tfrac43\pi R^3\Delta P + 4\pi R^2\sigma$, falls as $R$
grows, and the bubble grows by itself: the microcavitation. \emph{Proved here} from their
equation (1).
\end{enumerate}

Every link from the stretching to the microcavitation is read or follows from the review's own
equations.

\textbf{The chain as a seiche, held as theory.} Conceptually the chain is the same as a seiche:
energy aligned. A seiche is a standing wave in a basin, and its energy builds where the motion
arrives in phase. One part of this is read. Caupin and Herbert report that the acoustic experiments
reaching the most tension in water used aligned energy: Galloway a standing wave in a spherical
resonator, reaching $-20$ MPa in water degassed to $0.02$ percent saturation; Greenspan and
Tschiegg a standing wave focused in a steel cylinder, reaching $-16$ to $-21$ MPa at
$30\,^\circ$C; and their own experiment a few cycles of a 1 MHz wave tightly focused at the center
of a hemisphere, with an effective time of 45 nanoseconds. In each, the negative swing of the
aligned wave supplies the first link of the chain. That the chain and a seiche are the same in
concept is not in any source read, and no source on seiches is held. How long the compounding lasts in time, and what each link costs in a flow of equation
(1), are not in any source read. The transfer of angular momentum between bulk and boundary, and
the resemblance to quantum foam, are not in any source read either, and are not shown.

\section{The pattern left over}

Nine of the thirteen rows carry temperature. Viscosity, density, the speed of sound and bulk
viscosity each vary with it in the sources read. Heat is its balance. The level at which water
breaks moves with it in the review read, and the light of a collapsing bubble grows a hundredfold
as the water is cooled from $40$ to $0\,^\circ$C. The fluctuating stress is proportional to $k_BT$,
and gravity reaches the motion through $\rho(T)$. The four that do not are the two dependences on
pressure, of viscosity and of density, the molecules, and the external force. Pressure is the
second variable the rows share, and equation (1) holds pressure only as a gradient with no level,
Proposition 13.

Equation (1) has a place for every variable here except temperature, and Proposition 2 shows that a
viscous flow makes its own temperature gradient. The members equation (1) drops are therefore
mostly members the flow itself moves. That is the pattern left after the known parts are laid
side by side.

\textbf{Where the pattern leads.} The fluids of the world are transitive: every one of them carries
the members of the table, whatever its scale and however it is stirred. Along every route to an
unbounded value that this work has examined, one of those members acts first: the molecules,
thermal noise, the heat of shear, or the level at which the fluid breaks. A real fluid therefore
leaves the description of equation (1) before it could reach the value, and no singularity is
expected in observation, though other extreme objects in fluids are observed. Intuition I9 in the
log beside this workbook scores this: the first part as a hit for the routes examined, and the
absence of observation as open until a source on it is held.

\textbf{Status: theory.} The pattern is read off this table. It is not a law, it rests on the rows
entered so far, and a row added later can break it.

\textbf{The counter argument, held as theory.} In the author's words: ``Just like with every other
high energy reaction, the forces need to stay coupled until criticality.'' The forced construction
reaches its singular time by keeping two families of pulses aligned, so that their averaged
fluxes feed the core all the way down. A real fluid stays coupled that way only while it stays
liquid, its flow stays near laminar, and its viscosity and density stay what equation (1) holds
them at. Proposition 19 of the research paper proves that none of those can be decided in
equation (1): every solution of it, carried with a temperature, is compatible with a fluid past its
vapor point everywhere.

The other half is a scaling estimate. Propositions 15 and 18 set the temperature that shear holds
at a boundary by the square of the velocity difference across it, on the scale $\mu U^2/k$. Along
the forced construction the peak swirl grows as $\tau^{-1/2-h}$ (Duraiswami~\cite{src:Duraiswami-2026},
Table 7), so that scale grows as $\tau^{-1-2h}$; from the heating rate of the core and its radius,
the conducted rise grows as $\tau^{-1-3h}$. Both are without bound for the small positive $h$ of the
construction, while a liquid turns to vapor at a finite temperature. On this estimate a coupled
fluid lets go before the singular time, and in water the tension route of Proposition 17 lets go
sooner still, at 10 to 17 meters per second. The estimate is not checked against the computed
profile, and it is not a proof.

\section{What a laboratory would show}

Propositions 14 to 17 of the research paper put the first boundary water reaches at the cavity.
From atmospheric pressure at $20\,^\circ$C the axis of a turning column reaches the vapor pressure
at a swirl near $7.6$ meters per second, and shear warms a layer moving at the same speed by about
a hundredth of a kelvin. Duraiswami~\cite{src:Duraiswami-2026} proposes the apparatus: a swirl
chamber between porous walls driven toward collapse, with the onset of cavitation on the axis and
the loss of the steady state as what is observed.

\textbf{Uniform water.} The axis cavitates first, at the swirl where its pressure deficit reaches
$p_\infty - p_{\mathrm{sat}} + \Delta_n$. For a core of the Lamb--Oseen form that swirl is
$\sqrt{(p_\infty - p_{\mathrm{sat}} + \Delta_n)/(1.702\,\rho)}$. It rises with the pressure the
chamber is held at and falls as the gas dissolved in the water rises. No heating is measurable
before it: $0.06$ kelvin at $17$ meters per second.

\textbf{Water in layers.} The author's addition to I13 is scored a hit on the chemical difference.
Pimenova and Goldobin~\cite{src:Pimenova-and-Goldobin-2014} report water and n-heptane boiling at
their interface at $78.56\,^\circ$C while neither bulk boils, and Pfeiffer and
others~\cite{src:Pfeiffer-et-al-2023} report bubbles under tension nucleating mostly along the
interface between water and a perfluorocarbon that holds more gas. A column carrying such a layer
along its axis is expected to cavitate at a lower swirl than uniform water, and at the interface
first. A layer of warmer water lowers the swirl through $p_{\mathrm{sat}}$ alone: on the IAPWS
values, $7.46$ meters per second at $40\,^\circ$C, $6.97$ at $60\,^\circ$C and $5.71$ at
$80\,^\circ$C, against $7.63$ at $20\,^\circ$C, before any change in $\Delta_n$. A difference
of density alone has no source held, and nothing is predicted for it.

\textbf{What would count against it.} Vapor appearing in the bulk of uniform water away from the
axis before the axis cavitates; heating of a kelvin or more before the onset; or a column with such a layer
that cavitates at the same swirl as uniform water, or in the bulk before the interface.

\textbf{Status: theory.} The figures are computed from exact solutions and the IAPWS formulations,
and the expectations for layers are read off the two sources named. No such chamber has been
built or run for this work.
